\begin{abstract}
部分相关视频检索旨在根据只描述部分内容的文本查询，找出未经剪辑的视频。然而，简短查询与丰富视频内容之间固有的不对称，必然给检索过程带来\emph{不确定性}。在这一场景中，含糊的查询往往造成视频间\emph{语义歧义}；与此同时，视频内的\emph{稀疏时序监督}无法提供充分的匹配证据，进一步加剧了这一问题。

为此，我们提出分层证据学习框架 \modelname{}，通过汇聚不同粒度的跨模态证据，显式量化并建模不确定性。在\textbf{视频间}层面，我们将相似度分数视为证据支持，并以 Dirichlet 分布建模。基于所提出的三项判别原则，我们对查询进行细粒度识别，进而开展查询自适应的校准学习。在\textbf{视频内}层面，为积累更密集的证据，我们采用带自适应回收槽的灵活最优传输，构建查询与片段之间的软对齐，从而缓解时序监督稀疏，并抑制虚假局部响应。大量实验表明，\modelname{} 优于当前先进方法。代码发布于 \href{https://github.com/lijun2005/ICML26-Holmes}{ICML26-Holmes}。
\end{abstract}